% ============================================================
% Cheat sheet, argomento: il pattern Adapter (dispensa, capitolo 15)
% ============================================================
\argomento{Adapter}{Dispensa: capitolo 15. Pattern strutturale della GoF (sinonimo: Wrapper)}[Pattern Adapter]

\scheda{Carta d'identità}{
  \pintento{\textbf{Convertire l'interfaccia} di una classe esistente in quella che il client si aspetta, così che
    classi con interfacce incompatibili possano lavorare insieme. Come l'adattatore per le prese di corrente.}
  \ptesto{«riusare una classe (una libreria, un servizio, un microservizio) che esiste già», «la sua interfaccia non è
    quella che ci serve», «non possiamo modificarla» (non è nostra, o la usano altri), «unità di misura diverse»
    (libbre e chilogrammi), «formati diversi», «il nostro codice chiama \code{inEuro(\ldots)}, il servizio offre
    \code{getRate(\ldots)}».}
  \pcodice{una classe che \textbf{\code{implements}} l'interfaccia del client (Target) e ha un \textbf{campo del tipo della
    classe esistente} (Adaptee, di tipo \textbf{diverso} da Target), oppure la estende; i suoi metodi
    \textbf{traducono} parametri e risultati e chiamano un metodo con un altro nome o un'altra firma
    (\code{evaluate} chiama \code{calcolaIMC}).}
  \pdiagramma{\code{Client} $\to$ \code{«interface» Target}, realizzata da \code{Adapter}, che ha un'associazione verso
    \code{Adaptee} (adapter di oggetto) o ne è una sottoclasse (adapter di classe).}
  \pvarianti{adapter di \textbf{oggetto} (composizione: il più usato in Java, adatta anche le sottoclassi
    dell'Adaptee e lo nasconde) e adapter di \textbf{classe} (\code{extends Adaptee implements Target}:
    \vv{adapter-classe}). L'uso tipico: un'interfaccia decisa dal nostro progetto davanti a un servizio esterno.}
  \pconfondere{il \textbf{Decorator} (\vv{decorator-struttura}): avvolge un oggetto con la \textbf{stessa} interfaccia e
    aggiunge funzioni; l'Adapter avvolge un oggetto di un tipo \textbf{diverso} e ne cambia l'interfaccia. Il
    \textit{Facade}: semplifica un intero sottosistema, non adatta una classe (\vv{altri-gof}).}
}

\begin{struttura}{adapter-struttura}{La struttura: adapter di oggetto e di classe}
\begin{tikzpicture}
  \node[classe1=1.2cm] (c) at (0,0) {Client};
  \umlclass[2.2cm]{t}{(2.4,0)}{\textit{Target}}{}{+ request()}
  \umlclass[2.6cm]{a}{(2.4,-2.1)}{Adapter}{}{+ request()}
  \umlclass[2.8cm]{e}{(6.5,-2.1)}{Adaptee}{}{+ specificRequest()}
  \draw[dep] (c.east) -- (t.west);
  \draw[realizza] (a.north) -- (t.south);
  \draw[navig] (a.east) -- (e.west);
  \estO{(e.west)}{1}{adaptee}
  \node[notauml, anchor=north, text width=3.4cm] (n) at (2.4,-3.3) {adaptee.specificRequest()};
  \draw[tratt] (a.south) -- (n.north);
  \node[font=\tiny\bfseries, text=pcolore] at (6.0,0.3) {di oggetto (composizione)};
\end{tikzpicture}

\medskip
\begin{tikzpicture}
  \node[classe1=1.2cm] (c) at (0,0) {Client};
  \umlclass[2.2cm]{t}{(2.4,0)}{\textit{Target}}{}{+ request()}
  \umlclass[2.6cm]{a}{(2.4,-2.1)}{Adapter}{}{+ request()}
  \umlclass[2.8cm]{e}{(6.0,0)}{Adaptee}{}{+ specificRequest()}
  \draw[dep] (c.east) -- (t.west);
  \draw[realizza] (a.north) -- (t.south);
  \draw[gen] (a.east) -| (e.south);
  \node[font=\tiny\bfseries, text=pcolore] at (6.0,-1.3) {di classe (ereditarietà)};
\end{tikzpicture}
\tcblower
\textbf{Client}: vuole un servizio. \textbf{Target}: l'interfaccia che il client si aspetta. \textbf{Adaptee}: la
classe esistente, con un'interfaccia diversa. \textbf{Adapter}: implementa Target e traduce le chiamate verso
l'Adaptee, che \textbf{contiene} (sopra) o da cui \textbf{eredita} (sotto). Il client conosce solo Target: non sa che
l'Adaptee esiste.
\end{struttura}

\begin{programma}{adapter-esempio}{L'esempio del corso: l'indice di massa corporea}
\begin{lstlisting}[style=javacs, name=Main.java]
public class Main {
    public static void main(String[] args) {
        Health health = new Health();
        health.evaluateIMC(165, 70); // lb, pollici
        System.out.println(health.IMC);
    }
}
class Health {                   // il Client
    public double IMC;
    public double evaluateIMC(double weight,
                              double height) {
        Target adapter = new Adapter();
        IMC = adapter.evaluate(weight, height);
        return IMC;
    }
}
interface Target {               // ciò che serve
    double evaluate(double weight, double height);
}
class Adapter implements Target {
    private final Adaptee adaptee = new Adaptee();
    public double evaluate(double weight,
                           double height) {
        // libbre -> chilogrammi, pollici -> metri
        return adaptee.calcolaIMC(weight * 0.453,
                                  height * 0.0254);
    }
}
class Adaptee {                  // ciò che esiste
    public double calcolaIMC(double pesoKg,
                             double altezzaM) {
        return pesoKg / (altezzaM * altezzaM);
    }
}
\end{lstlisting}
\tcblower
\begin{tikzpicture}
  \umlclass[5.6cm]{h}{(0,0)}{Health}{+ IMC: double}{+ evaluateIMC(weight, height): double}
  \umlclass[5.6cm]{t}{(0,-2.0)}{\stereotipo{interface}{Target}}{}{+ evaluate(weight, height): double}
  \umlclass[5.6cm]{a}{(0,-4.0)}{Adapter}{}{+ evaluate(weight, height): double}
  \umlclass[5.6cm]{e}{(0,-6.0)}{Adaptee}{}{+ calcolaIMC(pesoKg, altezzaM): double}
  \draw[dep] (h.south) -- node[stereo, right]{«use»} (t.north);
  \draw[dep] (h.east) -- ++(0.6,0) |- node[stereo, right, pos=0.25]{«create»} (a.east);
  \draw[realizza] (a.north) -- (t.south);
  \draw[navig] (a.south) -- (e.north);
  \estN{(e.north)}{1}{adaptee}
\end{tikzpicture}

\medskip
\begin{lstlisting}[style=shellcs, name=Terminale]
$ javac Main.java
$ java Main
23.643873818359882
\end{lstlisting}
\regola{Tutta la «traduzione» sta nell'\code{Adapter}, in una riga: 165 libbre sono 74,7 kg, 70 pollici 1,78 m.
\code{Health} conosce solo \code{Target}. Nella slide della soluzione l'adapter \textbf{divide} per 0,453 e 0,0254, e
stampa un indice senza senso (0,000048): per passare dalle libbre ai chilogrammi si moltiplica. Nel codice del corso
l'Adaptee usa a sua volta una classe \code{IMC} con lo stesso calcolo.}
\end{programma}

\begin{caso}{adapter-classe}{La variante: adapter di classe}
\begin{lstlisting}[style=javacs]
interface Target {
    double evaluate(double weight, double height);
}
class Adaptee {
    public double calcolaIMC(double kg, double m) {
        return kg / (m * m);
    }
}
// eredita dall'Adaptee invece di contenerlo
class AdapterDiClasse extends Adaptee
                      implements Target {
    public double evaluate(double weight,
                           double height) {
        return calcolaIMC(weight * 0.453,
                          height * 0.0254);
    }
}
\end{lstlisting}
\tcblower
\begin{tikzpicture}
  \umlclass[4.5cm]{t}{(0,0)}{\stereotipo{interface}{Target}}{}{+ evaluate(weight, height): double}
  \umlclass[3.6cm]{e}{(4.35,0)}{Adaptee}{}{+ calcolaIMC(kg, m): double}
  \umlclass[4.5cm]{a}{(2.0,-2.5)}{AdapterDiClasse}{}{+ evaluate(weight, height): double}
  \draw[realizza] ($(a.north)+(-0.6,0)$) |- ++(0,0.4) -| (t.south);
  \draw[gen] ($(a.north)+(0.6,0)$) |- ++(0,0.4) -| (e.south);
\end{tikzpicture}
\regola{L'adapter \textbf{è} un Adaptee e chiama il metodo ereditato. Più compatto, ma legato a \textbf{una} classe
Adaptee, che espone tutta (chi ha un \code{AdapterDiClasse} può chiamare anche \code{calcolaIMC}); e in Java si
estende una classe sola. Per questo si preferisce l'adapter di oggetto.}
\end{caso}
